\documentclass[12pt,a4paper]{article} 

\renewcommand{\rmdefault}{ptm}   % Use Adobe Times

\usepackage{graphicx}
\usepackage{enumerate}
\usepackage{xtab}
\usepackage{pifont}
\usepackage{fancybox}
\RequirePackage[T1]{fontenc}
\RequirePackage[latin1]{inputenc}
\usepackage[portuges]{babel}
\RequirePackage{hyperref}
\hypersetup{
            pdfpagemode=UseThumbs,
            colorlinks=true,pdfhighlight=/P,linkbordercolor=1 1 1,
	        linkcolor=black,citecolor=black
	    } 

\usepackage[pdftex]{insdljs} 

\textwidth165mm
\textheight245mm
\oddsidemargin0mm
\evensidemargin0mm

\pagestyle{empty} 
\topmargin-10mm

\begin{insDLJS}[mysinus]{mysinus}{My Private sinus function}
function mysinus()
{
    var output = this.getField("Output.fx");
    output.value = Math.sin( Math.PI*x2.value/180 ) ;
    output.value = Math.round( output.value*1e10 )/1e10 ;
} 
\end{insDLJS}

\begin{insDLJS}[myarcsinus]{myarcsinus}{My Private arcsinus function}
function myarcsinus()
{
    var output = this.getField("Output.fx");
    if ( Math.abs( x2.value ) > 1 ) 
    { app.alert("O valor inserido é inválido!") }
    output.value = Math.asin( x2.value )*180/Math.PI ;
    output.value = Math.round( output.value*1e10 )/1e10 ;
} 
\end{insDLJS}

\begin{insDLJS}[mycosinus]{mycosinus}{My Private cosinus function}
function mycosinus()
{
    var output = this.getField("Output.fx");
    output.value = Math.cos( Math.PI*x2.value/180 ) ;
    output.value = Math.round( output.value*1e10 )/1e10 ;
} 
\end{insDLJS}

\begin{insDLJS}[myarccosinus]{myarccosinus}{My Private arccosinus function}
function myarccosinus()
{
    var output = this.getField("Output.fx");
    if ( Math.abs( x2.value ) > 1 ) 
    { app.alert("O valor inserido é inválido!") }
    output.value = Math.acos( x2.value )*180/Math.PI ;
    output.value = Math.round( output.value*1e10 )/1e10 ;
} 
\end{insDLJS}

\begin{insDLJS}[mytangens]{mytangens}{My Private tangens function}
function mytangens()
{
    var output = this.getField("Output.fx");
    output.value = Math.tan( Math.PI*x2.value/180 ) ;
    output.value = Math.round( output.value*1e10 )/1e10 ;
} 
\end{insDLJS}

\begin{insDLJS}[myarctangens]{myarctangens}{My Private arctangens function}
function myarctangens()
{
    var output = this.getField("Output.fx");
    output.value = Math.atan( x2.value )*180/Math.PI ;
    output.value = Math.round( output.value*1e10 )/1e10 ;
} 
\end{insDLJS}

\begin{insDLJS}[myexpx]{myexpx}{My Private exponential function}
function myexpx()
{
    var output = this.getField("Output.fx");
    output.value = Math.exp( x2.value ) ;
    output.value = Math.round( output.value*1e10 )/1e10 ;
} 
\end{insDLJS}

\begin{insDLJS}[mylnx]{mylnx}{My Private logarithm function}
function mylnx()
{
    var output = this.getField("Output.fx");
    output.value = Math.log( x2.value ) ;
    output.value = Math.round( output.value*1e10 )/1e10 ;
} 
\end{insDLJS}

\begin{insDLJS}[mysqrt]{mysqrt}{My Private square root function}
function mysqrt()
{
    var output = this.getField("Output.fx");
    output.value = Math.sqrt( x2.value ) ;
    output.value = Math.round( output.value*1e10 )/1e10 ;
} 
\end{insDLJS}

\begin{insDLJS}[mysquare]{mysquare}{My Private square function}
function mysquare()
{
    var output = this.getField("Output.fx");
    output.value = x2.value*x2.value ;
    output.value = Math.round( output.value*1e10 )/1e10 ;
} 
\end{insDLJS}

\begin{insDLJS}[equadratica]{equadratica}{equadratica}
function equadratica() 
{ 
var x1 = this.getField("Output.x1")
var x2 = this.getField("Output.x2")
var a = InputA.value  
var b = InputB.value  
var c = InputC.value  
var d = b*b - 4*a*c ; 
if ( a == 0 ) { 
app.alert("O coeficiente a deve ser diferente de zero!")
               } else {
if ( d >= 0 ) { 	       
x1.value = ( -b + Math.sqrt( d ) )/( 2*a ) ;
x2.value = ( -b - Math.sqrt( d ) )/( 2*a ) ;
             }	else { 
app.alert("As raízes desta equação não são reais...")	      
	              }
	              }
x1.value = Math.round( x1.value*1e10 )/1e10 ;
x2.value = Math.round( x2.value*1e10 )/1e10 ;
}
\end{insDLJS}

%%%%%%%%%%%%%%%%%%%%%%%%%%%%%%%%%%%%%%%%%%%%%%%%%%%%%%%%%%%%%%%%%%%%%%%%%%%%%
%%%%%%%%%%%%%%%%%%%%%%%%%%%%%%%%%%%%%%%%%%%%%%%%%%%%%%%%%%%%%%%%%%%%%%%%%%%%%

\begin{document}

%%%%%%%%%%%%%%%%%%%%%%%%%%%%%%%%%%%%%%%%%%%%%%%%%%%%%%%%%%%%%%%%%%%%%%%%%%%%%
%%%%%%%%%%%%%%%%%%%%%%%%%%%%%%%%%%%%%%%%%%%%%%%%%%%%%%%%%%%%%%%%%%%%%%%%%%%%%

\begin{Form}

\section{Funções mais comuns}
\hrule
\vskip3mm
\hrule
\vskip5mm 
\begin{center}
Insira o valor de $x$ \\ 
e clique nos botões para efectuar o respectivo cálculo:
\end{center}
\vskip15mm
\begin{center}
\begin{tabular}{ccc}
$x$    & = & 
\TextField[width=60mm,name=Input.x2,borderstyle=I,value={0},align=1,
validate={x2=this.getField("Input.x2");}]{ } \\ 
\ \\
$f(x)$ & = &
\TextField[width=60mm,name=Output.fx,borderstyle=I,value={0},align=1,
validate={fx=this.getField("Output.fx");}]{ } 
\end{tabular}\\
\vspace{10mm}
\begin{tabular}{ccc}
\PushButton[name=sinx,color=0 0 1,bordercolor=0 0 0,borderstyle=B,borderwidth=1.5,onclick={mysinus();}]
{
\begin{minipage}{30mm}
\begin{center}
$\sin x_{}$
\end{center}
\end{minipage}
} & \parbox{30mm}{} & 
\PushButton[name=cosx,color=0 0 1,bordercolor=0 0 0,borderstyle=B,borderwidth=1.5,onclick={mycosinus();}]
{
\begin{minipage}{30mm}
\begin{center} 
$\cos x_{}$
\end{center}
\end{minipage}
} \\ 
\ \\ 
\PushButton[name=tanx,color=0 0 1,bordercolor=0 0 0,borderstyle=B,borderwidth=1.5,onclick={mytangens();}]
{
\begin{minipage}{30mm}
\begin{center}
$\tan x_{}$
\end{center}
\end{minipage}
} & \ & \ \\ 
\ \\
\PushButton[name=arcsinx,color=0 0 1,bordercolor=0 0 0,borderstyle=B,borderwidth=1.5,onclick={myarcsinus();}]
{
\begin{minipage}{30mm}
\begin{center}
$\arcsin{x}$ 
\end{center}
\end{minipage}
} & \parbox{30mm}{} & 
\PushButton[name=arccosx,color=0 0 1,bordercolor=0 0 0,borderstyle=B,borderwidth=1.5,onclick={myarccosinus();}]
{
\begin{minipage}{30mm}
\begin{center}
$\arccos{x}$ 
\end{center}
\end{minipage}
} \\ 
\ \\
\PushButton[name=arctanx,color=0 0 1,bordercolor=0 0 0,borderstyle=B,borderwidth=1.5,onclick={myarctangens();}]
{
\begin{minipage}{30mm}
\begin{center}
$\arctan{x}$ 
\end{center}
\end{minipage}
} & \ & \ \\
\ \\
\PushButton[name=expx,color=0 0 1,bordercolor=0 0 0,borderstyle=B,borderwidth=1.5,onclick={myexpx();}]
{
\begin{minipage}{30mm}
\begin{center}
$e_{}^{x}$ 
\end{center}
\end{minipage}
} & \parbox{30mm}{} & 
\PushButton[name=lnx,color=0 0 1,bordercolor=0 0 0,borderstyle=B,borderwidth=1.5,onclick={mylnx();}]
{
\begin{minipage}{30mm}
\begin{center}
$\ln x_{}$
\end{center}
\end{minipage}
} \\ 
\ \\
\PushButton[name=sqrt,color=0 0 1,bordercolor=0 0 0,borderstyle=B,borderwidth=1.5,onclick={mysqrt();}]
{
\begin{minipage}{30mm}
\begin{center}
$\sqrt{x_{}}$ 
\end{center}
\end{minipage}
} & \parbox{30mm}{} & 
\PushButton[name=square,color=0 0 1,bordercolor=0 0 0,borderstyle=B,borderwidth=1.5,onclick={mysquare();}]
{
\begin{minipage}{30mm}
\begin{center}
$x_{}^2$
\end{center}
\end{minipage}
} \\ 
\end{tabular}
\end{center} 
\vskip5mm  
\hrule
\vskip3mm
\hrule

\clearpage 

\section{Raízes da equação quadrática}

\hrule
\vskip3mm
\hrule
\vskip5mm 
\begin{center}
Insira os valores de $a$, $b$ e $c$\\ 
e clique no botão para calcular as raízes da equação quadrática:
\end{center}
\vskip15mm
\begin{center}
\shadowbox{
\begin{tabular}{c}
$ax^2 + bx + c = 0$ \\
\ \\ 
\ \\ 
$x_{1,2} = \displaystyle{ \frac{-b\pm\sqrt{b^2-4ac}}{2a} }$
\end{tabular}
           }
\end{center}
\vskip10mm
\begin{center}
\begin{tabular}{ccc}
$a$ = \TextField[width=30mm,name=InputA,borderstyle=I,value={},align=1,
validate={InputA=this.getField("InputA");}]{ } & 
$b$ = \TextField[width=30mm,name=InputB,borderstyle=I,value={},align=1,
validate={InputB=this.getField("InputB");}]{ } & 
$c$ = \TextField[width=30mm,name=InputC,borderstyle=I,value={},align=1,
validate={InputC=this.getField("InputC");}]{ } 
\ \\ 
\ \\ 
\ & 
\PushButton[name=pbequadratica,borderstyle=B,borderwidth=1.5,onclick={equadratica();}]
{\hskip2mm Mostrar as raizes\hskip2mm} & \ \\
\ \\ 
$x_1$ = \TextField[width=30mm,name=Output.x1,borderstyle=I,value={},align=1,
validate={x1=this.getField("Output.x1");}]{ } \ & \ & 
$x_2$ = \TextField[width=30mm,name=Output.x2,borderstyle=I,value={},align=1,
validate={x2=this.getField("Output.x2");}]{ } 
\end{tabular}
\end{center}
\vskip5mm  
\hrule
\vskip3mm
\hrule

\end{Form}

%%%%%%%%%%%%%%%%%%%%%%%%%%%%%%%%%%%%%%%%%%%%%%%%%%%%%%%%%%%%%%%%%%%%%%%%%%%%%
%%%%%%%%%%%%%%%%%%%%%%%%%%%%%%%%%%%%%%%%%%%%%%%%%%%%%%%%%%%%%%%%%%%%%%%%%%%%%

\end{document}

%%%%%%%%%%%%%%%%%%%%%%%%%%%%%%%%%%%%%%%%%%%%%%%%%%%%%%%%%%%%%%%%%%%%%%%%%%%%%
%%%%%%%%%%%%%%%%%%%%%%%%%%%%%%%%%%%%%%%%%%%%%%%%%%%%%%%%%%%%%%%%%%%%%%%%%%%%%
